\subsection*{NT-L16. Zero multiplicity forces logarithmic mean-square growth}
For a fixed primitive unitary Hecke function, use the inverse coefficients
and summatory function $\mathcal M_\eta$ of NT-L15, and put
\[
 I_\eta(X)=\int_1^X|\mathcal M_\eta(x)/x|^2dx.
\]
If $\rho=1/2+i\gamma$ is a zero of multiplicity $m\ge1$, write
$c_\rho=L^{(m)}(\rho,\eta)/m!$. Then, without assuming GRH,
\[
 \limsup_{X\to\infty}\frac{I_\eta(X)}{(\log X)^{2m-1}}
       \ge\frac1{(2m-1)((m-1)!)^2\,|\rho c_\rho|^2}>0.
                                                        \tag{NT-L16a}
\]
More strongly, the reciprocal leading-coefficient terms may be summed
over \emph{all distinct critical-line zeros of exact multiplicity $m$}:
\[
 \limsup_{X\to\infty}\frac{I_\eta(X)}{(\log X)^{2m-1}}
 \ge\frac1{(2m-1)((m-1)!)^2}
   \sum_{\substack{\Re\rho=1/2\\ \operatorname{ord}_\rho L=m}}
            \frac1{|\rho c_\rho|^2}.                    \tag{NT-L16a'}
\]
This includes the possibility that the sum and the limsup are infinite.
Consequently, for every $0\le a<2m-1$,
\[
 \limsup_{X\to\infty}I_\eta(X)/(\log X)^a=+\infty.
                                                        \tag{NT-L16b}
\]
If instead there is a zero with real part greater than $1/2$, then
(NT-L16b) holds for every finite real $a\ge0$.
The critical-line location of the zero in (NT-L16a) is an explicit
hypothesis; $M$ supplies it for every nontrivial zero, but is not
used in the proof of this implication.

\begin{proof}
Set $F(y)=e^{-y/2}\mathcal M_\eta(e^y)$ and
$A(Y)=\int_0^Y|F(y)|^2dy=I_\eta(e^Y)$.
Let $b=2m-1$ and
$B=\limsup_{Y\to\infty}A(Y)/Y^b$.
If $B=+\infty$, both asserted lower bounds are immediate.
Otherwise, for every $\epsilon>0$,
$A(Y)\le(B+\epsilon)Y^b+C_\epsilon$.
As in NT-L15, Cauchy--Schwarz with an exponential weight gives a
holomorphic Laplace transform on $\Re w>0$, and partial summation
in its initial absolute-convergence half-plane gives
\[
 H(w)=\int_0^\infty F(y)e^{-wy}dy
      =\frac1{(w+1/2)L(w+1/2,\eta)}.
                                                        \tag{NT-L16c}
\]
The identity extends meromorphically through $\Re w>0$.
Use again the inner product
\[
 \langle f,g\rangle_\sigma=\sigma\int_0^\infty
       f(y)\overline{g(y)}e^{-\sigma y}dy.
\]
Integration by parts yields
\[
 \begin{split}
 \|F\|_\sigma^2
 &=\sigma^2\int_0^\infty A(y)e^{-\sigma y}dy\\
 &\le (B+\epsilon)\Gamma(b+1)\sigma^{1-b}
                               +C_\epsilon\sigma.
 \end{split}                                           \tag{NT-L16d}
\]
Put $k=m-1$ and use the normalized polynomial waves
\[
 u_{\gamma,\sigma}(y)=
       \frac{\sigma^k y^k e^{i\gamma y}}{\sqrt{(2k)!}},
 \qquad v_\sigma(y)=\sigma^k F(y).
\]
Their Gram matrix has entries
\[
 \langle u_{\gamma,\sigma},u_{\nu,\sigma}\rangle_\sigma
       =\left(\frac{\sigma}{\sigma-i(\gamma-\nu)}\right)^{2k+1}.
                                                        \tag{NT-L16e}
\]
In particular each vector has norm one, and for any fixed finite set
of distinct ordinates its Gram matrix tends to the identity as
$\sigma\downarrow0$. Polynomial growth of $A$ permits differentiation
under the Laplace integral on every compact subset of $\Re w>0$,
by weighted Cauchy--Schwarz. Hence
\[
 \langle v_\sigma,u_{\gamma,\sigma}\rangle_\sigma
 =\frac{\sigma^{2k+1}}{\sqrt{(2k)!}}
                 (-1)^k H^{(k)}(\sigma+i\gamma).
\]
At a zero $\rho=1/2+i\gamma$ of exact order $m=k+1$,
the meromorphic identity (NT-L16c) gives
$H(w)=(\rho c_\rho)^{-1}(w-i\gamma)^{-(k+1)}$
plus terms of smaller pole order. Differentiating this finite
principal part yields
\[
 \lim_{\sigma\downarrow0}
 \langle v_\sigma,u_{\gamma,\sigma}\rangle_\sigma
       =\frac{\sqrt{(2k)!}}{k!\,\rho c_\rho}.           \tag{NT-L16f}
\]
For a fixed finite set $S$ of the specified zeros, finite-dimensional
orthogonal projection onto the span of these waves gives the Gram
matrix version of Bessel's inequality. By (NT-L16e), letting
$\sigma\downarrow0$ in this inequality and using (NT-L16d) gives
\[
 \frac{(2k)!}{(k!)^2}\sum_{\rho\in S}|\rho c_\rho|^{-2}
   \le \limsup_{\sigma\downarrow0}\|v_\sigma\|_\sigma^2
   \le (B+\epsilon)(2k+1)!.
\]
Let $\epsilon\downarrow0$ and take the supremum over finite $S$.
Since the terms are nonnegative, this proves (NT-L16a') and
therefore (NT-L16a). For $m=1$ it is exactly the Abel--Bessel
normalization of NT-L15. The one-frequency model
$F(y)=a y^{m-1}e^{i\gamma y}/(m-1)!$ has
$A(Y)=|a|^2Y^{2m-1}/((2m-1)((m-1)!)^2)$, so this constant is
sharp for the single-pole Laplace argument; no attainment by an
arithmetic Hecke function is asserted.
The positive limsup in (NT-L16a) gives (NT-L16b) along an unbounded sequence
when $a<2m-1$.

Finally suppose a zero has real part greater than $1/2$.
A finite limsup in (NT-L16b), for any fixed $a\ge0$, would imply
$A(Y)\le C(1+Y^a)$. The same weighted Cauchy--Schwarz argument
makes $H$ holomorphic on $\Re w>0$. Equation (NT-L16c) would then
exclude every zero of $L$ in $\Re s>1/2$, a contradiction.
This proves the final assertion.
\end{proof}
This is a direct multiplicity-sensitive version of the weak-Mertens
Laplace argument. It is a necessary analytic growth condition, not
an asymptotic formula for $I_\eta(X)$.

\subsection*{NT-L17. Actual Kummer fields obstruct universal weak Mertens and full-jet positivity}
Let $r\ge3$ be an odd prime and let $p\ne r$ be a rational prime. Put
\[
 K=\mathbb Q(\zeta_r),\qquad E=\mathbb Q(\zeta_r,p^{1/r}).
\]
There is a nontrivial primitive finite-order Hecke character $\theta$
of $K$ whose entire finite-part $L$-function $H(s)=L_K(s,\theta)$
satisfies the all-prime factorization
\[
                  \zeta_E(s)=\zeta_K(s)H(s)^{r-1}.     \tag{NT-L17a}
\]
The function $H$ has infinitely many nontrivial zeros. Hence $\zeta_E$
has infinitely many nontrivial zeros of multiplicity at least $r-1$,
\emph{unconditionally}.

Let $\mathcal M_E(x)=\sum_{N\mathfrak a\le x}\mu_E(\mathfrak a)$
and $I_E(X)=\int_1^X|\mathcal M_E(x)/x|^2dx$. Then
\[
            \limsup_{X\to\infty}I_E(X)/\log X=+\infty.
                                                        \tag{NT-L17b}
\]
Thus a universal weak-Mertens hypothesis for all Dedekind zeta
functions is false. This statement makes no assertion about an
individual Dirichlet function's weak-Mertens conjecture.

If RH for $\zeta_E$ is additionally assumed, choose any zero $\rho$
of $H$. It is then on the critical line. With
$m=\operatorname{ord}_\rho\zeta_E\ge r-1$, NT-L16 gives the explicit
nonzero lower bound
\[
 \limsup_{X\to\infty}\frac{I_E(X)}{(\log X)^{2m-1}}
 \ge
 \frac1{(2m-1)((m-1)!)^2\,
             |\rho\,\zeta_E^{(m)}(\rho)/m!|^2}.
                                                        \tag{NT-L17c}
\]
In particular $\limsup I_E(X)/(\log X)^{2r-3}>0$; it is infinite
if the chosen multiplicity exceeds $r-1$. If instead $\zeta_E$
has an off-critical zero, its functional equation supplies a
right-half-strip zero, and every polynomial-logarithmic mean bound
fails by NT-L16. This gives the same positive-limsup conclusion
unconditionally by a dichotomy, while (NT-L17c) retains its stated
critical-zero/RH hypothesis.

There is also a concrete consequence for MT-M20: the canonical
\emph{full-jet} quotient for $\zeta_E$ has no positive-definite
continuous inner product for which every convolution extends boundedly
and every $C_{h*h^\sharp}$ is positive. Under $M$, the Hilbert traces
in MT-M19 are nevertheless nonnegative. Trace positivity and a faithful
positive operator realization are different assertions here.

\begin{proof}
The elementary Galois calculation in CX-ALG-024 gives
$[E:K]=r$ and
$\operatorname{Gal}(E/\mathbb Q)=\mathbb F_r\rtimes\mathbb F_r^\times$.
For clarity, $p$ is unramified in the cyclotomic field $K$, and
$X^r-p$ is Eisenstein at every prime above $p$, proving the degree.
The translations send $p^{1/r}$ to $\zeta_r^a p^{1/r}$.
The automorphism fixing $p^{1/r}$ and sending $\zeta_r$ to
$\zeta_r^b$ conjugates a translation by $a$ to one by $ba$.
Thus the $r-1$ nontrivial characters of
$\operatorname{Gal}(E/K)=\mathbb F_r$ form a single orbit under
$\operatorname{Gal}(K/\mathbb Q)$.

Class field theory identifies these characters with primitive
finite-order Hecke characters $\theta^a$ of $K$.
Conjugating ideals by an automorphism of $K/\mathbb Q$ preserves
absolute norms and carries Frobenius and inertia by conjugation.
Consequently it permutes every finite Euler factor, including the
ramified factors, and all $L_K(s,\theta^a)$ with $a\ne0$ are the
same function $H(s)$. The unconditional all-prime abelian factorization
proved within NT-H11 therefore gives (NT-L17a).
Equivalently, CX-ALG-024's permutation-character identity gives
$H=\zeta_{\mathbb Q(p^{1/r})}/\zeta_{\mathbb Q}$.
No general Artin-holomorphy conjecture is being invoked: $H$ is an
actual nontrivial Hecke $L$-function, so it is entire by the classical
Hecke theorem (Ramakrishnan--Valenza, Theorem7-19, printed p.279).

Greni\'e--Molteni, AppendixA TheoremA.1, author p.26, applies
unconditionally to primitive finite-order Hecke characters. Its
zero count has main term
$([K:\mathbb Q]/\pi)T\log T+O_{K,\theta}(T)$.
Hence $H$ has infinitely many nontrivial zeros in $0<\Re s<1$.
At each such zero $\zeta_K$ is holomorphic, because its only pole
is at $1$. Equation (NT-L17a) therefore gives multiplicity at
least $r-1$ in $\zeta_E$.
If the limsup in (NT-L17b) were finite, NT-L15 applied to the
trivial primitive Hecke character over $E$ would make all these
zeros simple, a contradiction. This proves (NT-L17b).
The quantitative claims follow exactly from NT-L16 and the
functional equation of $\zeta_E$, as described above.
Finally MT-M20 proves that a faithful positive completion with the
stated convolution properties requires simple zeros. The actual
multiple zeros just proved exclude it. MT-M19 supplies the separate
trace-positive conclusion under RH for $\zeta_E$, which $M$ supplies
conditionally.
\end{proof}
For example $r=3$, $p=2$ gives the degree-six $S_3$ field
$E=\mathbb Q(\sqrt[3]{2},\sqrt{-3})$ and at least double nontrivial
zeros. The argument given here concerns these explicitly monomial
Galois fields; this proof alone does not address every nonabelian
Galois field.
Geometrically this is the positive multiplicity of a single
rank-one Hecke realization in the zeta realization of an affine
finite Galois cover. It is compatible with GRH and with the
unconditional adapted-block synthesis of MT-M18, but prevents
a faithful positive Hilbert realization of the full jet algebra,
even when its traces are positive under $M$.
